\documentclass[11pt,letterpaper]{article}
\usepackage[T1]{fontenc}
\usepackage[utf8]{inputenc}
\usepackage{lmodern,microtype}
\usepackage[margin=0.78in,headheight=15pt]{geometry}
\usepackage{amsmath,amssymb,booktabs,longtable,tabularx,array,enumitem,xcolor,hyperref,fancyhdr,lastpage,needspace}
\definecolor{EngineBlue}{HTML}{183B56}
\definecolor{EngineTeal}{HTML}{117A83}
\definecolor{EngineGray}{HTML}{596671}
\definecolor{EngineLight}{HTML}{EEF4F7}
\hypersetup{colorlinks=true,urlcolor=EngineTeal,linkcolor=EngineBlue}
\setlength{\parindent}{0pt}
\setlength{\parskip}{5pt}
\emergencystretch=3em
\setlist{nosep,leftmargin=*}
\pagestyle{fancy}\fancyhf{}
\fancyhead[L]{\small\sffamily Game Engine Architecture | Engineering Suite}
\fancyhead[R]{\small\sffamily \DocType}
\fancyfoot[L]{\scriptsize Source-grounded engineering curriculum | v1.0}
\fancyfoot[R]{\scriptsize \thepage/\pageref{LastPage}}
\newcommand{\DocType}{Reference}
\newcommand{\SuiteHeading}[3]{%
\renewcommand{\DocType}{#2}%
{\small\bfseries\sffamily\color{EngineTeal}#2\hfill #3}\par\vspace{6pt}
{\LARGE\bfseries\color{EngineBlue}#1}\par\vspace{4pt}
{\color{EngineTeal}\rule{\linewidth}{1pt}}\par\vspace{7pt}}
\newcommand{\LabEvidence}[1]{\par\textbf{Evidence required:} #1\par}
\newcommand{\WriteLines}[1]{\par\begingroup\setlength{\parskip}{9pt}\count255=0\loop\ifnum\count255<#1\noindent\rule{\linewidth}{0.35pt}\par\advance\count255 by1\repeat\endgroup\smallskip}
\newcommand{\WritingBox}[1]{\noindent\fbox{\parbox[t][#1][t]{0.97\linewidth}{\mbox{}}}\par\smallskip}
\newcommand{\SourceNote}{\footnotesize\textit{Primary reading: Jason Gregory, \emph{Game Engine Architecture}, 3rd ed. (CRC Press, 2018). The narrative, exercises, laboratory designs and scoring instruments in this suite are independently authored teaching materials, not a reproduction of the textbook.}\normalsize\par}
\newcommand{\Criterion}[2]{\textbf{#1.} #2\par}
\newcommand{\FormField}[2]{\Needspace{3\baselineskip}\textbf{#1}\par\textcolor{EngineGray}{\footnotesize #2}\par\WriteLines{2}}
\newcommand{\BlankField}[1]{\Needspace{3\baselineskip}\textbf{#1}\par\WriteLines{2}}
\newcolumntype{Y}{>{\raggedright\arraybackslash}X}

\begin{document}
\SuiteHeading{Laboratory 04: Generational Resource Handles}{Implementation lab}{C++}
\SourceNote
\section{Reading and entry gate}
Book chapters: 6, 7, 16. Language: C++. Associated laboratory code: code/cpp/generation\_handles.cpp. Prior gates: correct local toolchain, a deterministic unit test and documented hardware/platform.

\section{Learning goals and implementation contract}
\begin{enumerate}
\item Implement slot/generation IDs; reuse slots with incremented generation.
\item Ensure stale IDs resolve to no object after destruction.
\item State generation wrap and over-concurrent-read limitations.
\end{enumerate}
\section{Execution}
From this suite root, build C++ using CMake in code/cpp; run Python 3 scripts in code/python. Prefer a debug build with sanitizers for correctness and a separate optimized build for measurements. No commercial engine or copyrighted assets are required.

\section{Required negative and verification tests}
\begin{itemize}
\item Negative test for stale handle.
\item Duplicate erase should be safe.
\item Resource owner and thread affinity are explicit.
\end{itemize}
\section{Deliverables}
\begin{itemize}
\item Minimal reproducible repository state with commands and commit identifier.
\item Input record, expected results and actual results with the failing case included.
\item One-page interface/lifetime diagram and alternative decision.
\item Benchmark CSV (where relevant) and interpretation of confounding variables.
\end{itemize}
\section{Lab review questions}
\begin{enumerate}
\item Which invariant would detect a subtle implementation bug?
\item What state must be synchronized or kept immutable?
\item How does behavior change when content size or frame rate doubles?
\item Which results are implementation evidence versus theoretical expectation?
\end{enumerate}
\section{Assessment}
\small\begin{longtable}{|p{0.313\linewidth}|p{0.313\linewidth}|p{0.313\linewidth}|}
\hline \textbf{Criterion} & \textbf{Points} & \textbf{Evidence}\\ \hline\endhead
Design contract & 20 & Interface and lifetime\\ \hline
Working implementation & 30 & Passing tests\\ \hline
Adversarial validation & 25 & Negative tests\\ \hline
Reproducibility and metrics & 15 & Raw evidence\\ \hline
Trade-off explanation & 10 & Written decision\\ \hline
\end{longtable}\normalsize

\end{document}
